\documentclass[11pt,a4paper]{report}
\usepackage{amsmath,amssymb}
\usepackage{booktabs}
\usepackage{hyperref}
\usepackage{xcolor}
\usepackage{geometry}
\usepackage{enumitem}
\usepackage{longtable}
\usepackage{graphicx}
\usepackage{quoting}
\usepackage{multirow}
\usepackage{array}
\geometry{margin=0.9in}
\setlength{\parskip}{0.4em}
\setlength{\parindent}{0em}

\definecolor{q}{RGB}{90,90,90}
\definecolor{finding}{RGB}{0,80,140}
\newcommand{\cq}[1]{\textcolor{q}{\textit{``#1''}}}
\newcommand{\finding}[1]{\textcolor{finding}{\textbf{Finding:}} #1}

\title{Three Walls to 1.0\\[0.5em]\large MRI Forensics on Five Causal Bank Checkpoints\\Instance 7 (Heinrich + Chronohorn), April 11, 2026}

\author{asuramaya \and Claude Opus 4.6 (1M context, Instance 7)\\[0.3em]\small \url{https://github.com/asuramaya/chronohorn}}

\date{April 2026}

\begin{document}
\maketitle
\tableofcontents

%==============================================================
\part{The Paper}
%==============================================================

\chapter{Introduction}

The causal bank architecture achieves 1.80 bits per byte (bpb) on FineWeb English text at 12.4M parameters. English has an estimated entropy of $\sim$1.0 bpb. That leaves 0.80 bpb the model cannot predict.

This paper uses Heinrich, a model forensics tool, to perform complete MRI captures on five causal bank checkpoints trained by Chronohorn. The MRI captures substrate states, expert routing decisions, per-band logit contributions, embedding vectors, and all learned weights for every token in the vocabulary. We then run 200 sequences of length 512 through the best model with full internal capture to measure per-position loss, routing dynamics, substrate--local path balance, and mode utilization across the sequence.

The result is a precise decomposition of the 0.80 bpb gap into three independent walls:

\begin{enumerate}
\item \textbf{Raw byte tax} (62\% of gap, 0.46 bpb): the 1024-token vocabulary forces 0.9\% of tokens into byte-fallback mode where they score 54.7 bpb. This is a tokenizer problem, not a model problem.
\item \textbf{Architectural ceiling} (35\% of gap, 0.26 bpb): even at optimal operating conditions (positions 4--63, non-raw tokens), the model reaches 1.26 bpb. The remaining 0.26 bpb represents computations that exponential moving averages and 8-token local windows cannot perform.
\item \textbf{Position degradation} (4\% of gap, 0.03 bpb): substrate mode saturation causes a mild loss increase at late positions. The readout already compensates for most of the EMA growth.
\end{enumerate}

Along the way, twelve mechanistic questions are answered using the MRI data. Each finding is independently verifiable from the captured arrays.


\chapter{Experimental Setup}

\section{Checkpoints}

Five checkpoints were transferred from the Chronohorn fleet to an Apple Silicon workstation for MRI capture:

\begin{table}[h]
\centering
\begin{tabular}{lrrrrl}
\toprule
\textbf{Checkpoint} & \textbf{Scale} & \textbf{Params} & \textbf{bpb} & \textbf{Steps} & \textbf{Notes} \\
\midrule
frozen-proj-asym8420-s12 & 12 & 14.6M & 1.882 & 10k & asym experts, frozen proj \\
learned-proj-asym8420-s12 & 12 & 14.8M & 1.890 & 10k & asym experts, learned proj \\
push-exp2-ema2000-s12-50k & 12 & 12.4M & 1.801 & 50k & best model, hl$_{\max}$=2000 \\
push-exp2-ema1000-s12-50k & 12 & 12.4M & 1.802 & 50k & hl$_{\max}$=1000 \\
scale16-asym8420-w8 & 16 & 22.2M & 1.873 & 10k & larger scale \\
\bottomrule
\end{tabular}
\caption{Five causal bank checkpoints under analysis.}
\label{tab:checkpoints}
\end{table}

All models share:
\begin{itemize}[nosep]
\item Vocabulary: 1024 SentencePiece BPE tokens (256 byte tokens + 768 merges)
\item Local window: 8 tokens, scaling factor 0.25
\item Readout: routed squared-ReLU experts
\item Training data: FineWeb-10B tokenized with sp1024
\item Substrate: frozen random input projection, EMA decay modes
\end{itemize}

The ``asym8420'' models use asymmetric expert allocation: 8 experts in band 0 (fast modes, half-life 1.5--8), 4 in band 1 (medium, hl 8--39), 2 in band 2 (slow, hl 39--197), and 0 in band 3 (slowest, hl 199--1000), where band 3 is replaced by a static \texttt{nn.Linear(embed\_dim, vocab\_size)} that takes only the embedding as input.

The ``push'' models use symmetric 2-expert allocation across all 4 bands with no static band.

\section{MRI Capture}

Heinrich's MRI tool captures per-token impulse responses via the decepticon backend. For each token $t \in [0, 1023]$, the model receives a single-token input from zero initial state and records:

\begin{itemize}[nosep]
\item \textbf{substrate.npy}: $[1024, n_{\text{modes}}]$ float16 --- EMA state at readout
\item \textbf{routing.npy}: $[1024, n_{\text{experts}}]$ float16 --- expert routing weights
\item \textbf{band\_logits.npy}: $[1024, 4, 1024]$ float16 --- per-band logit contributions
\item \textbf{embedding.npy}: $[1024, d_{\text{embed}}]$ float16 --- input embeddings
\item \textbf{half\_lives.npy}: $[n_{\text{modes}}]$ float32 --- mode decay half-lives
\item \textbf{weights/}: all learned parameters as individual .npy files
\end{itemize}

Capture time: 2 seconds per model on Apple M3 Max. All five MRIs completed in under 15 seconds.

The decepticon loader required a fix to support asymmetric expert allocation. The original \texttt{load\_checkpoint()} read \texttt{linear\_readout\_num\_experts: 2} uniformly from the result JSON, which doesn't capture per-band expert counts. We added \texttt{\_infer\_band\_experts()}, which inspects the checkpoint's \texttt{state\_dict} keys to reconstruct the asymmetric allocation (e.g., 8 experts in band 0 from the presence of \texttt{\_band\_readouts\_0\_experts\_in\_\{0..7\}}), and patches the config before model construction.

\section{Sequence-Level Analysis}

For questions requiring sequential dynamics (per-position loss, routing stability, substrate growth), we loaded the push-ema2000-50k model directly and ran \texttt{forward\_captured()} on 200 sequences of length 512 drawn from the FineWeb validation set (62M tokens, sp1024-tokenized). Two outlier tokens ($>$ 1023) were removed from the raw data before reshaping.

Loss was computed as cross-entropy in nats, then converted to bpb using the tokenizer's measured tokens-per-byte ratio of 0.4105.


\chapter{The Three Walls}

\section{Wall 1: Raw Byte Tax (0.46 bpb, 62\% of gap)}

\finding{0.9\% of target tokens are byte-fallback tokens scoring 54.7 bpb. They contribute 0.46 bpb to the overall average. When excluded, the model operates at 1.285 bpb.}

The sp1024 vocabulary contains 256 byte tokens (IDs 0--255) and 768 BPE merges (IDs 256--1023). In practice, only 93 of the 256 byte tokens appear in the validation data---the remaining 163 are dead. The byte tokens that do appear represent UTF-8 fragments that the tokenizer couldn't merge: \texttt{0x00} and \texttt{0x01} (39\% of raw bytes, likely sequence boundary artifacts), high-byte characters like \texttt{0xe6} (6.3\%, CJK fragments), and \texttt{0xc6} (4.6\%, accented character fragments).

\begin{table}[h]
\centering
\begin{tabular}{lrrr}
\toprule
\textbf{Target class} & \textbf{bpb} & \textbf{\% of tokens} & \textbf{Contribution} \\
\midrule
Raw bytes (0--255) & 54.67 & 0.9\% & 0.469 bpb \\
Early merges (256--355) & 1.19 & 35.1\% & 0.419 bpb \\
Mid merges (356--555) & 1.46 & 20.7\% & 0.302 bpb \\
Late merges (556--755) & 1.58 & 10.1\% & 0.159 bpb \\
Final merges (756--1023) & 1.18 & 33.3\% & 0.394 bpb \\
\midrule
\textbf{Overall} & \textbf{1.74} & \textbf{100\%} & \textbf{1.743 bpb} \\
\bottomrule
\end{tabular}
\caption{Loss by BPE merge rank. Raw bytes dominate the gap despite being $<$1\% of tokens.}
\label{tab:merge_rank}
\end{table}

Raw byte frequency is uniform across sequence positions (0.43--0.63\%), ruling out position-dependent concentration. The raw bytes are uniformly unpredictable because the 1024-vocab tokenizer has insufficient coverage to merge them---this is a vocabulary capacity problem, not a modeling failure.

The BPE merge rank gradient is mild: early merges (common sub-words) at 1.19 bpb vs.\ late merges (rarer composites) at 1.58 bpb, a 33\% gap. Final merges (756--1023) perform as well as early merges (1.18 bpb) because they represent common multi-character sequences comprising 33\% of all tokens.

\subsection{Loss as Input}

The merge rank also affects prediction of the \emph{following} token:

\begin{table}[h]
\centering
\begin{tabular}{lr}
\toprule
\textbf{After...} & \textbf{Following token bpb} \\
\midrule
Raw bytes & 58.66 \\
Early merges & 1.36 \\
Mid merges & 1.23 \\
Late merges & 1.16 \\
Final merges & 1.17 \\
\bottomrule
\end{tabular}
\caption{Loss of the token following each merge rank class. Raw byte inputs corrupt prediction of the next token.}
\label{tab:following_loss}
\end{table}

After a raw byte, the next token is also unpredictable (58.7 bpb), suggesting that raw bytes cluster in sequences of unmergeable fragments. After mid and late merges, the following token is \emph{easier} to predict (1.16--1.23 bpb), consistent with these tokens appearing in structured, predictable contexts.


\section{Wall 2: Architectural Ceiling (0.26 bpb, 35\% of gap)}

\finding{At optimal operating conditions (positions 4--63, non-raw tokens), loss is 1.258 bpb. The remaining 0.26 bpb above English entropy is what exponential moving averages cannot compute.}

\subsection{Per-Position Loss}

\begin{table}[h]
\centering
\begin{tabular}{lrr}
\toprule
\textbf{Position range} & \textbf{All tokens (bpb)} & \textbf{Non-raw only (bpb)} \\
\midrule
0--3 (cold start) & 1.573 & 1.548 \\
4--15 (warming) & 1.291 & 1.246 \\
16--63 (optimal) & 1.307 & 1.261 \\
64--127 & 1.322 & 1.237 \\
128--255 & 1.527 & 1.254 \\
256--510 & 2.064 & 1.314 \\
\midrule
\textbf{Overall} & \textbf{1.743} & \textbf{1.285} \\
\bottomrule
\end{tabular}
\caption{Per-position loss. The dramatic late-position degradation (2.06 bpb) largely disappears when raw bytes are excluded (1.31 bpb), revealing that substrate saturation contributes only +0.07 bpb of real degradation.}
\label{tab:position_loss}
\end{table}

Excluding raw bytes reveals that the model maintains near-constant performance from position 4 through position 256. The apparent catastrophic late-position degradation at 2.06 bpb was an artifact: raw bytes at late positions score $\sim$172 bpb (far worse than the global raw byte average of 55 bpb), inflating the mean.

The non-raw loss at position 4--63 is 1.258 bpb. This is the model's floor---the best it can achieve with the substrate fully warmed and context available. The 0.258 bpb gap between this floor and English entropy (1.0 bpb) represents computations the architecture fundamentally cannot perform.

\subsection{What the EMA Cannot Compute}

Cross-architecture comparison confirms the representation gap. We computed linear centered kernel alignment (CKA) between the causal bank's substrate states and Qwen 0.5B's hidden states at every layer:

\begin{table}[h]
\centering
\begin{tabular}{lr}
\toprule
\textbf{Comparison} & \textbf{CKA} \\
\midrule
CB substrate vs.\ Qwen L0 & 0.083 \\
CB substrate vs.\ Qwen L6 & 0.010 \\
CB substrate vs.\ Qwen L12 & 0.010 \\
CB substrate vs.\ Qwen L18 & 0.010 \\
CB substrate vs.\ Qwen L23 & 0.023 \\
CB embedding vs.\ Qwen embedding & 0.209 \\
CB final logits vs.\ Qwen L23 & 0.015 \\
\bottomrule
\end{tabular}
\caption{CKA between causal bank (push-ema2000-50k) and Qwen 0.5B. Different tokenizers (sp1024 vs 151k) make this a geometry comparison. CKA is essentially zero: the two architectures learn completely different representations.}
\label{tab:cka}
\end{table}

CKA $\approx 0$ at every layer. Even the embedding comparison (0.209) is low, and it's the highest because both encode token identity. By Qwen L6, alignment is 0.01---the representations have fully diverged.

This means the 0.26 bpb gap isn't information the causal bank is ``missing''---it's information that requires a fundamentally different computation. Three capabilities that attention provides and EMA does not:

\begin{enumerate}
\item \textbf{Selective recall}: ``the cat that the dog chased was...'' requires attending back to ``cat'', not averaging over all prior tokens. The EMA integrates everything with exponential decay. It cannot skip, select, or compare.
\item \textbf{Compositional binding}: ``A is red, B is blue'' requires separate variable slots, not a weighted mixture. The substrate represents one mixed state, not a set of bindings.
\item \textbf{Long-range syntax}: bracket matching, agreement across clauses. The 8-token local window handles immediate $n$-gram prediction but cannot reach further.
\end{enumerate}


\section{Wall 3: Position Degradation (0.03 bpb, 4\% of gap)}

\finding{Substrate mode amplification is real (18$\times$ ramp for slow modes) but the readout compensates. Non-raw loss increases only 0.07 bpb from best to worst positions.}

The EMA theory predicts that slow modes accumulate signal over their half-life window. Mode utilization measurements from 50 sequences of length 512 confirm this precisely:

\begin{table}[h]
\centering
\begin{tabular}{lrrrrl}
\toprule
\textbf{Band} & \textbf{Half-life} & \textbf{Early (0--16)} & \textbf{Late (256--512)} & \textbf{Ramp} & \\
\midrule
Band 0 & 2--9 & 0.135 & 0.157 & 1.2$\times$ & flat \\
Band 1 & 9--54 & 0.189 & 0.440 & 2.3$\times$ & mild \\
Band 2 & 55--330 & 0.243 & 2.341 & 9.6$\times$ & strong \\
Band 3 & 334--2001 & 0.254 & 4.585 & 18.0$\times$ & explosive \\
\bottomrule
\end{tabular}
\caption{Mean absolute substrate activation by band and position range. The ramp ratio tracks half-life as predicted: fast modes are position-invariant, slow modes accumulate.}
\label{tab:mode_ramp}
\end{table}

No dead modes were found: all 516 modes fire above the noise floor at every position. The five most position-varying modes are all in band 3 (half-lives 981--1789). Mode 501 (hl=1644.5) goes from 0.09 at position 0 to 24.0 at position 512---a 263$\times$ dynamic range.

Despite this explosive growth, the band logit magnitudes in the readout output tell a different story about how the model handles it:

\begin{table}[h]
\centering
\begin{tabular}{lrrrr}
\toprule
& \textbf{Pos 0--8} & \textbf{Pos 8--64} & \textbf{Pos 64--256} & \textbf{Pos 256--512} \\
\midrule
Band 0 logit L2 & 124 & 138 & 142 & 140 \\
Band 1 logit L2 & 82 & 105 & 134 & 139 \\
Band 2 logit L2 & 29 & 25 & 177 & 718 \\
Band 3 logit L2 & 27 & 57 & 1550 & 7695 \\
\midrule
Substrate total & 262 & 324 & 2004 & 8693 \\
Local ($\times$0.25) & 333 & 259 & 255 & 257 \\
\textbf{Ratio (sub/local)} & \textbf{0.8$\times$} & \textbf{1.3$\times$} & \textbf{7.8$\times$} & \textbf{33.8$\times$} \\
\bottomrule
\end{tabular}
\caption{Band logit L2 norms and substrate--local path balance by position. The substrate dominates after position 64 and overwhelms the local path by 34$\times$ at late positions. Band 3 grows from 27 to 7695.}
\label{tab:substrate_local}
\end{table}

The substrate path produces 34$\times$ larger logits than the local path at late positions, but the non-raw loss only increases by 0.07 bpb. The readout experts have learned to extract signal from growing substrate magnitudes. The substrate saturation is a dynamic range challenge the model has mostly solved internally, not an architecture limitation.


\chapter{Twelve Questions from MRI}

\section{Q1: Learned vs.\ Frozen Projection Geometry}

\cq{Does the learned projection show band-aligned structure, or did it just add noise?}

\finding{The learned projection IS structured but makes things 0.008 bpb WORSE. Conflicting-gradient hypothesis confirmed.}

The frozen-proj (1.882 bpb) and learned-proj (1.890 bpb) checkpoints differ only in substrate mode: \texttt{frozen} vs.\ \texttt{learnable\_mixing}. PCA on the substrate states reveals strikingly different structure:

\begin{table}[h]
\centering
\begin{tabular}{lrr}
\toprule
\textbf{Metric} & \textbf{Frozen projection} & \textbf{Learned projection} \\
\midrule
PC1 variance & 3.6\% & \textbf{7.3\%} \\
PCs for 50\% & 37 & \textbf{21} \\
PCs for 95\% & 190 & \textbf{150} \\
Effective dimensionality & 153.8 & \textbf{100.9} \\
\bottomrule
\end{tabular}
\caption{PCA comparison of substrate states. The learned projection concentrates variance into fewer dimensions.}
\label{tab:pca_comparison}
\end{table}

The frozen projection distributes PC loadings uniformly across bands (25\% per band per PC), as expected for a random matrix. The learned projection concentrates on bands 0 and 2 (40--43\% loading) while suppressing band 3 to 9--11\%:

\begin{verbatim}
  Frozen PC1 (3.6%): B0=25% B1=32% B2=20% B3=23%
  Learned PC1 (7.3%): B0=26% B1=22% B2=43% B3=9%
  Learned PC2 (4.9%): B0=41% B1=23% B2=26% B3=10%
\end{verbatim}

Cross-projection PCs are nearly orthogonal (max cosine 0.22). The learned projection found a \emph{completely different subspace}---not a refinement of the random one. Per-mode cosine between the two projections averages 0.64 (36\% divergence).

The interpretation: gradient descent found band-aligned structure in the projection (concentrating on fast and medium-slow modes, ignoring the slowest band), but this structure created conflicting objectives between bands. The different bands need different projection geometries, and a single shared projection cannot satisfy all of them simultaneously. The random projection avoids this conflict by treating all bands equally, which turns out to be 0.008 bpb better.


\section{Q2: 10k vs.\ 50k Expert Routing}

\cq{Did the within-band experts specialize further, or just get better at the same specialization?}

\finding{MRI impulse response shows 100\% expert-0 collapse (zero-state artifact). Sequence context reveals band-dependent routing dynamics.}

\subsection{MRI Routing (Single-Token Impulse)}

All three models route 100\% of tokens to expert 0 in the MRI capture. This is because the MRI feeds single tokens from zero initial EMA state---the router's bias term dominates when there is no substrate context. The routing array is zero for the second expert across all 1024 tokens. This is a property of the capture mode, not the model.

\subsection{Sequence Routing}

With hooks into each band's \texttt{RoutedSquaredReLUReadout}, we captured per-band expert routing on real FineWeb sequences:

\begin{table}[h]
\centering
\begin{tabular}{lrrr|rrr}
\toprule
& \multicolumn{3}{c}{\textbf{push-ema2000}} & \multicolumn{3}{c}{\textbf{push-ema1000}} \\
\textbf{Band} & \textbf{E0\%} & \textbf{E1\%} & \textbf{Switch} & \textbf{E0\%} & \textbf{E1\%} & \textbf{Switch} \\
\midrule
B0 (hl 2--9) & 48 & 52 & 34\% & 44 & 56 & 36\% \\
B1 (hl 9--54) & 73 & 27 & \textbf{1.8\%} & 90 & 10 & \textbf{3.8\%} \\
B2 (hl 55--330) & 54 & 46 & 37\% & 82 & 18 & 14\% \\
B3 (hl 334--2001) & 49 & 51 & 41\% & 40 & 60 & 39\% \\
\bottomrule
\end{tabular}
\caption{Per-band expert routing in sequence context. Band 1 routes at the sequence level (1.8--3.8\% switch rate). Other bands switch at the token level.}
\label{tab:routing_sequence}
\end{table}

Three regimes emerge:

\begin{enumerate}
\item \textbf{Band 1 (medium timescale): sequence-level routing.} Expert 0 wins 73--90\% of tokens with only 1.8--3.8\% switch rate. The medium-decay modes integrate enough context ($\sim$50 positions) to make a stable expert choice early and keep it.
\item \textbf{Bands 0, 2, 3: token-level routing.} 34--41\% switch rate means the expert changes every 2--3 tokens. These bands respond to local token content, not sequence-level features.
\item \textbf{ema1000 collapses more.} Band 1: 90/10 (vs.\ 73/27 for ema2000). Band 2: 82/18 (vs.\ 54/46). Shorter half-life range provides less diverse substrate features, so the router has less to differentiate on.
\end{enumerate}

Band 1's routing also shows position dependence: it starts at 22--24\% expert 0 for the first 16 positions, then jumps to 74--94\% for the remainder. The expert needs $\sim$16 positions of context to decide.


\section{Q3: Where the 0.80 bpb Hides}

\cq{Where in the sequence is the model worst?}

\finding{The loss is U-shaped across positions. Cold start at positions 0--3, optimal at 4--127, mild degradation at 128+. But 62\% of the total gap comes from raw byte tokens, not position effects.}

This question produced the key decomposition of the paper, presented in Chapter 3. The position curve for non-raw tokens is remarkably flat:

\begin{verbatim}
  Pos   4-15:  1.246 bpb (non-raw)
  Pos  16-64:  1.261 bpb
  Pos  64-128: 1.237 bpb  <-- model's best
  Pos 128-256: 1.254 bpb
  Pos 256-511: 1.314 bpb
\end{verbatim}

The sweet spot at position 64--128 (1.237 bpb) likely corresponds to where the medium-timescale EMA modes have accumulated useful context without the slow modes having saturated. This aligns with the Q9 mode utilization data: band 1 (hl 9--54) reaches steady state around position 50--100.

Additional breakdown by token frequency:
\begin{itemize}[nosep]
\item Rare tokens (Q1 by frequency): 3.37 bpb --- a 3.5$\times$ gap vs.\ common
\item Common tokens (Q4 by frequency): 0.97 bpb --- \emph{below} English entropy
\item After punctuation: 1.48 bpb (0.28 bpb easier than average)
\end{itemize}


\section{Q4: Scale 16 Routing---8 Experts All Differentiate}

\cq{At 8 experts, are they all differentiated, or do some collapse into pairs?}

\finding{All 8 experts are fully differentiated. The specialization ceiling has not been hit.}

The scale16-asym8420 checkpoint has 8 experts in band 0. We analyzed both the router weight vectors and the expert input projection weights:

\begin{table}[h]
\centering
\begin{tabular}{lrrrr}
\toprule
\textbf{Model} & \textbf{Band} & \textbf{$n$ experts} & \textbf{Router cos (mean)} & \textbf{Expert\_in cos (mean)} \\
\midrule
Scale 16 & B0 & 8 & $-0.137$ & 0.024 \\
Scale 16 & B1 & 4 & $-0.231$ & 0.031 \\
Scale 16 & B2 & 2 & $-0.920$ & 0.096 \\
\midrule
Scale 12 & B0 & 8 & $-0.136$ & 0.024 \\
Scale 12 & B1 & 4 & $-0.240$ & 0.028 \\
Scale 12 & B2 & 2 & $-0.824$ & 0.091 \\
\bottomrule
\end{tabular}
\caption{Expert differentiation by band and scale. Negative router cosines indicate repulsion (good). Expert\_in cosines near zero indicate orthogonal learned representations (full differentiation).}
\label{tab:expert_diff}
\end{table}

The router cosine similarity is negative (repulsive) in all cases, indicating the router learned to partition the input space. The expert\_in weight cosines of 0.024 (band 0, both scales) mean the expert projection matrices are essentially orthogonal---each expert learned a completely different feature transformation.

The full 8$\times$8 cosine matrix for scale 16 band 0 routers shows no pair structure:

\begin{verbatim}
  [ 1.000 -0.123  0.030 -0.028 -0.109 -0.263 -0.363 -0.321]
  [-0.123  1.000 -0.005 -0.094 -0.121 -0.191 -0.346 -0.147]
  [ 0.030 -0.005  1.000 -0.070 -0.048 -0.224 -0.317 -0.430]
  [-0.028 -0.094 -0.070  1.000 -0.183 -0.207 -0.131 -0.280]
  [-0.109 -0.121 -0.048 -0.183  1.000 -0.213  0.089 -0.198]
  [-0.263 -0.191 -0.224 -0.207 -0.213  1.000 -0.065  0.238]
  [-0.363 -0.346 -0.317 -0.131  0.089 -0.065  1.000  0.292]
  [-0.321 -0.147 -0.430 -0.280 -0.198  0.238  0.292  1.000]
\end{verbatim}

The highest pair similarity is (E6, E7) at 0.292 in the router, but their expert\_in weights have only 0.039 cosine---they route similarly on some inputs but compute different things. More experts in band 0 may still help.


\section{Q5: The Static Band 3 Weight Matrix}

\cq{What did band 3's \texttt{nn.Linear} learn? Unigram? Classifier? Gate?}

\finding{Band 3 is an input-dependent bigram table: $W \cdot \text{embed}[\text{input}] + b$. The $W$ component is 5.85$\times$ stronger than the bias.}

In the asymmetric models, band 3 has no substrate input---it's a bare \texttt{nn.Linear(384, 1024)} that maps the current token's embedding directly to logits. SVD decomposition of the weight matrix:

\begin{itemize}[nosep]
\item Top singular value: 40.4, effective rank: 110 (out of 384)
\item Rank for 50\% variance: 32; for 80\%: 113; for 95\%: 228
\item Bias entropy: 9.99 bits (maximum is 10.0)---the bias is nearly uniform, not a unigram
\item Weight logit std: 0.574 vs.\ bias logit std: 0.098, ratio 5.85$\times$
\end{itemize}

The near-uniform bias (9.99/10.0 bits entropy) rules out a unigram distribution. The $W$ matrix dominates: it projects each input embedding to a token-specific logit pattern. This is a \emph{learned bigram table}---given the current token, predict the next token via a linear transformation of the embedding, with no temporal context.

Band 3's contribution is modest: L2 norm 17.2, compared to band 0's 123.0 (14\%). It provides a context-free prior that the substrate-dependent bands can override.


\section{Q6: Cross-Band Logit Interference}

\cq{Do the bands agree or fight?}

\finding{Asymmetric models: bands cooperate (cos 0.45--0.92). Symmetric push models: band 3 actively suppresses band 0's predictions 48--67\% of the time.}

\begin{table}[h]
\centering
\begin{tabular}{lrrr}
\toprule
\textbf{Model} & \textbf{cos(B0,B1)} & \textbf{cos(B0,B3)} & \textbf{B3 suppresses B0 top-1} \\
\midrule
frozen-proj (asym, 10k) & 0.920 & 0.449 & 0.6\% \\
push-ema2000 (sym, 50k) & 0.862 & 0.724 & \textbf{47.6\%} \\
push-ema1000 (sym, 50k) & 0.860 & 0.718 & \textbf{67.2\%} \\
scale16 (asym, 10k) & 0.923 & 0.499 & 0.7\% \\
\bottomrule
\end{tabular}
\caption{Cross-band logit interference. Asymmetric models agree; symmetric models fight.}
\label{tab:cross_band}
\end{table}

The asymmetric architecture forces agreement: band 3's static bigram lookup has no way to ``know'' that band 0 is wrong about a particular token. The symmetric push models, where band 3 has a full routed expert readout with substrate access, learned to use slow bands as \textbf{correction signals}. When band 0's fast modes make an initial prediction based on recent tokens, band 3's slow modes (which have accumulated longer context) override 48--67\% of the time.

Top-1 agreement rates with the final prediction:
\begin{itemize}[nosep]
\item frozen-proj: B0=41\%, B3=43\% (band 3 matches final more often than band 0)
\item push-ema2000: B0=31\%, B3=3.6\% (band 3 rarely determines the final prediction alone)
\end{itemize}

In the push model, band 3 works by \emph{pushing the wrong answer down}, not by \emph{pushing its own answer up}. Its top-1 agrees with the final only 3.6\% of the time, yet it contributes through suppression.


\section{Q7: Expert Routing Stability}

Covered in Q2. Summary: band 1 routes at the sequence level (1.8\% switch rate), bands 0/2/3 route at the token level (34--41\% switch rate). The routing timescale of band 1 matches its EMA timescale: medium-decay modes integrate enough context to make stable decisions.


\section{Q8: Substrate vs.\ Local Path}

\cq{When does the substrate dominate? When does the local path dominate?}

\finding{Local path dominates at positions 0--8 (0.8$\times$). Substrate takes over by position 8--64 (1.3$\times$) and overwhelms at 256--512 (33.8$\times$). The substrate is not decoration.}

Covered in detail in Section 3.3 (Table \ref{tab:substrate_local}). The local path provides a stable baseline of $\sim$257 L2 norm regardless of position. The substrate path grows from 262 to 8693. At cold start, the local window is the only source of prediction. By position 64, the substrate has accumulated enough signal to dominate.

The crossover point is position 8--64, where the substrate-to-local ratio passes 1.0. This aligns with the loss curve: the model reaches optimal performance at position 4--64, exactly where the substrate begins contributing without yet saturating.


\section{Q9: Mode Utilization Heatmap}

\cq{Do fast modes fire at every position? Do slow modes ramp up gradually?}

\finding{EMA theory confirmed exactly. Fast modes are position-invariant (1.2$\times$ ramp). Slow modes ramp 18$\times$. Zero dead modes.}

Covered in Section 3.3 (Table \ref{tab:mode_ramp}). The key additional finding: the five most position-varying modes are all in band 3 with half-lives $>$ 981. Mode 501 (hl=1644.5) has a 263$\times$ dynamic range between position 0 and position 512. This mode represents the longest-memory signal in the substrate: it integrates over approximately 1644 positions, far beyond the 512-position sequence length. Even at position 512, it hasn't reached steady state.


\section{Q10: Input Projection Null Space}

\cq{Which modes are in the null space? Do they correlate with band 3's behavior?}

\finding{The null space is uniformly distributed. The random projection has no systematic bias toward any band.}

The frozen projection maps $[384 \to 516]$: 384 embedding dimensions to 516 modes. With rank 384, at least 132 modes must be in the null space. The singular values are tightly clustered (0.152 to 2.159), consistent with a random Gaussian matrix. Mode norms have coefficient of variation 0.036 (3.6\%)---essentially uniform.

Leverage scores (how well each mode is represented by the projection) show no band dependence:

\begin{table}[h]
\centering
\begin{tabular}{lrr}
\toprule
\textbf{Band} & \textbf{Leverage mean} & \textbf{Half-life range} \\
\midrule
B0 (fast) & 0.748 & 1.5--7.6 \\
B1 (medium) & 0.745 & 7.6--38.5 \\
B2 (slow) & 0.742 & 39--196 \\
B3 (slowest) & 0.742 & 199--1000 \\
\bottomrule
\end{tabular}
\caption{Leverage scores by band. Correlation with half-life: $-0.018$ (zero). The null space falls uniformly across bands.}
\label{tab:leverage}
\end{table}

14 modes fall below mean $-$ 2$\sigma$ leverage, distributed as [2, 5, 4, 3] across bands---no concentration in any band. Band 3's routing and readout behavior is not explained by projection quality; the random projection treats all modes identically.

The learned projection diverged significantly from frozen (global cosine 0.60, per-mode cosine 0.64$\pm$0.22), confirming that gradient descent moved the projection substantially---but this movement hurt performance (Q1). The random projection is empirically near-optimal.


\section{Q11: Token-Level Loss by BPE Merge Rank}

\cq{Is the model's bottleneck rare-token prediction?}

\finding{No. The bottleneck is raw byte tokens (0.9\% of data, 54.7 bpb). Among BPE merges, the loss gradient from common to rare is only 33\%.}

Covered in Section 3.1 (Table \ref{tab:merge_rank}). The key insight: final merges (token IDs 756--1023) at 1.18 bpb are as easy as early merges (1.19 bpb), despite representing higher merge ranks. This is because final merges comprise 33\% of all tokens---they're common multi-character units like frequent words. The BPE merge rank is not a good proxy for token rarity in this vocabulary.


\section{Q12: Causal Bank vs.\ Transformer CKA}

\cq{Do the two architectures capture similar information?}

\finding{CKA $\approx 0$ at every layer. They solve completely different optimization problems.}

Covered in Section 3.2 (Table \ref{tab:cka}). The caveat: different tokenizers (sp1024 vs.\ Qwen 151k) prevent token-level alignment, so this is a representation geometry comparison over the first 1024 token indices in each model's vocabulary. But the result is unambiguous: the similarity is at the noise floor.

The embedding CKA of 0.209 is the only non-trivial value, reflecting that both models must encode token identity. After even one layer of processing (Qwen L0, CKA=0.083), the representations diverge, and by L6 they reach 0.010.


\chapter{The Path to 1.0}

\section{Step 1: Vocabulary Expansion (Expected: 0.35 bpb)}

The raw byte tax is 62\% of the gap. Going from 1024 to 4096 tokens eliminates most byte-fallback scenarios:

\begin{itemize}[nosep]
\item 93 of 256 byte tokens appear in validation data; most are UTF-8 fragments
\item A 4096-token vocabulary would merge the most common fragments into proper tokens
\item The model already handles high-merge-rank tokens well (1.18 bpb for final merges)
\item Estimated gain: 0.30--0.40 bpb (eliminating the 0.47 bpb raw byte contribution minus residual rare-token loss)
\end{itemize}

This is pure infrastructure: no architecture change, no gradient interaction, no hyperparameter search. Retokenize the data, retrain.

The risk is minimal. The readout scales as $O(V \cdot d_{\text{hidden}})$ per expert; going from 1024 to 4096 quadruples the readout parameters (from $\sim$800k to $\sim$3.2M per expert), which is still small relative to the substrate and embedding.

\section{Step 2: Sparse Attention (Expected: 0.15 bpb)}

The architectural ceiling is 0.26 bpb. The EMA substrate cannot perform selective recall, compositional binding, or long-range syntax. But full self-attention is $O(n^2)$ and defeats the purpose of the causal bank architecture.

Candidate mechanisms that preserve O($n$) complexity:

\begin{enumerate}
\item \textbf{Cross-attention every $K$ positions}: at position $t$, attend to the last $K$ substrate snapshots. The substrate provides compressed context; attention selects from it. Cost: $O(n/K \cdot K^2) = O(nK)$.
\item \textbf{Linear attention}: approximate the attention kernel with a feature map $\phi(Q)\phi(K)^T$. Runs as a recurrence alongside the EMA. Cost: $O(nd^2)$.
\item \textbf{Gated retrieval}: learn a gate that decides when to ``snapshot'' the substrate state into a fixed-size memory bank. The readout can then attend to the bank. Cost: $O(nM)$ for $M$ memory slots.
\end{enumerate}

Expected gain: 0.10--0.20 bpb, capturing the selective recall component of the architectural ceiling. Compositional binding ($\sim$0.05 bpb) and deep syntax ($\sim$0.05 bpb) may require more sophisticated mechanisms.

\section{Step 3: Scale and Training (Expected: 0.15 bpb)}

The remaining gap after vocabulary and attention:

\begin{itemize}[nosep]
\item \textbf{More modes}: currently 516 modes from a 384-dim embedding (overcomplete by 132). Scale to 1024 modes from a 512-dim embedding for better coverage.
\item \textbf{Deeper readout}: the current readout is single-layer per expert. A 2-layer readout can extract more from the substrate features.
\item \textbf{Longer training}: 50k steps on FineWeb-10B. The push models show no sign of converging (overfit is $-1.8\%$, meaning test loss is still improving faster than train loss). 200k steps should capture another 0.05 bpb.
\item \textbf{Fix routing collapse}: band 1 collapses to 90/10 in ema1000. Higher \texttt{balance\_coeff} or per-band load balancing could recover 0.02 bpb.
\end{itemize}

\section{What Not To Do}

The MRI findings specifically rule out several plausible-sounding improvements:

\begin{enumerate}
\item \textbf{Don't learn the input projection.} Q1 shows the learned projection finds band-aligned structure but performs 0.008 bpb worse. The conflicting-gradient hypothesis is confirmed: bands need different projection geometries, and a single shared learned projection cannot satisfy all of them. Keep the random projection.

\item \textbf{Don't normalize substrate states.} Q9 shows the substrate grows 18$\times$ at late positions, and Q3 shows this contributes only 0.03 bpb of degradation. The readout already compensates. Layer normalization would destroy the magnitude information that the readout has learned to use.

\item \textbf{Don't add more experts without adding attention.} Q4 shows 8 experts fully differentiate (cos 0.024). The experts are already orthogonal---the problem isn't expert capacity, it's the information reaching them. Without attention, more experts just partition the same EMA signal more finely.

\item \textbf{Don't expand modes beyond $\sim$512.} Q10 shows the random projection already has 132 null-space modes at 516 modes from 384 embedding dims. Adding modes without expanding the embedding just adds noise the projection can't distinguish.

\item \textbf{Don't optimize the tokenizer at fixed vocab size.} Q11 shows the loss gradient across merge ranks is only 33\% (1.19 to 1.58 bpb). The problem isn't merge order---it's vocabulary size. Reordering merges within 1024 tokens won't eliminate the byte-fallback tokens.
\end{enumerate}


\chapter{Supplementary Data}

\section{Checkpoint Architecture Details}

All models use the \texttt{causal\_bank\_torch} preset with:

\begin{verbatim}
  linear_impl:          fft
  linear_readout_kind:  routed_sqrelu_experts
  linear_readout_depth: 1
  balance_coeff:        0.01
  mix_mode:             additive
  init_policy:          chronohorn_v1
  oscillatory_frac:     0.0
  local_window:         8
  block_mixing_ratio:   0.25
  share_embedding:      false
\end{verbatim}

\begin{table}[h]
\centering
\begin{tabular}{lrrrrr}
\toprule
& \textbf{frozen-proj} & \textbf{learned-proj} & \textbf{push-2000} & \textbf{push-1000} & \textbf{scale16} \\
\midrule
embedding\_dim & 384 & 384 & 384 & 384 & 512 \\
linear\_modes & 516 & 516 & 516 & 516 & 512 \\
linear\_hidden & (767,) & (767,) & (767,) & (767,) & (1023,) \\
local\_hidden & (1536,) & (1536,) & (1536,) & (1536,) & (2048,) \\
hl\_min & 1.5 & 1.5 & 1.5 & 1.5 & 1.5 \\
hl\_max & 1000 & 1000 & 2000 & 1000 & 1000 \\
readout\_bands & 4 & 4 & 4 & 4 & 4 \\
band\_experts & (8,4,2,0) & (8,4,2,0) & (2,2,2,2) & (2,2,2,2) & (8,4,2,0) \\
substrate\_mode & frozen & learnable & frozen & frozen & frozen \\
tensors & 80 & 81 & 56 & 56 & 80 \\
train\_steps & 10k & 10k & 50k & 50k & 10k \\
\bottomrule
\end{tabular}
\caption{Full architecture parameters for all five checkpoints.}
\label{tab:full_arch}
\end{table}

\section{Training Results}

\begin{table}[h]
\centering
\begin{tabular}{lrrrr}
\toprule
& \textbf{Train bpt} & \textbf{Test bpt} & \textbf{Test bpb} & \textbf{Overfit \%} \\
\midrule
frozen-proj & 4.598 & 4.584 & 1.882 & $-0.29$ \\
learned-proj & --- & --- & 1.890 & --- \\
push-ema2000 & 4.466 & 4.386 & 1.801 & $-1.80$ \\
push-ema1000 & --- & --- & 1.802 & --- \\
scale16 & 4.574 & 4.562 & 1.873 & $-0.26$ \\
\bottomrule
\end{tabular}
\caption{Training metrics. Negative overfit means test loss is lower than train loss (no overfitting).}
\label{tab:training}
\end{table}

\section{MRI Capture Metadata}

All captures performed on Apple M3 Max, April 11, 2026. Heinrich v0.5, decepticon backend.

\begin{verbatim}
  Mode:     raw (single-token impulse from zero state)
  Tokens:   1024 (full vocabulary)
  Storage:  substrate [1024, 516/512] float16
            routing   [1024, 2] float16
            band_logits [1024, 4, 1024] float16
            embedding [1024, 384/512] float16
            half_lives [516/512] float32
            weights/  (all parameters as .npy)
  Time:     2 seconds per model
\end{verbatim}

Sequence-level analysis: 200 sequences $\times$ 512 positions $=$ 102,400 tokens from FineWeb validation set. Forward pass with full capture (\texttt{forward\_captured()}) recording substrate states, band logits, and local logits at every position.


\chapter{Conclusions}

The 0.80 bpb gap between the causal bank (1.80) and English entropy (1.0) decomposes into three independent walls:

\begin{enumerate}
\item \textbf{Vocabulary} (0.46 bpb, 62\%): fix by expanding from 1024 to 4096 tokens.
\item \textbf{Architecture} (0.26 bpb, 35\%): fix by adding sparse attention to the substrate.
\item \textbf{Saturation} (0.03 bpb, 4\%): already handled by the readout; no fix needed.
\end{enumerate}

The single highest-leverage change is vocabulary expansion. It requires no architecture modification, no gradient research, and removes 62\% of the gap. The second change---sparse attention---is the research frontier: finding an O($n$) mechanism that captures selective recall without destroying the causal bank's throughput advantage.

The MRI forensics approach---capturing complete internal states and analyzing them offline---proved essential. The naive analysis (overall loss by position) showed a ``substrate saturation catastrophe'' at late positions (2.06 bpb). Decomposing by token type revealed this was almost entirely a vocabulary artifact (raw bytes scoring 172 bpb at late positions). Without the MRI data to separate these effects, we would have spent engineering effort on a 4\%-of-gap problem while ignoring the 62\%-of-gap solution sitting in the tokenizer.


\end{document}
